\chapter{美国证券从业资格考试}

本章覆盖 FINRA 旗下的核心证券从业资格考试：SIE、Series 7、Series 57、Series 65、Series 24。
这些考试是进入美国证券行业的关键门槛，也是量化金融从业者理解\textbf{监管框架与业务规则}的重要参考。
与其他章节不同，本章的重点是\textbf{可计算}：每个知识领域都配有公式与数值例题。

\begin{warningbox}[考试参数具有时效性]
本章列出的题量、时长、通过分数与费用来自编写时的 FINRA 公开大纲。
\textbf{FINRA 会定期调整考试结构、费用与通过标准}，报考前务必以
\texttt{finra.org} 的最新《Content Outline》为准。
\end{warningbox}

\begin{table}[htbp]
\centering
\caption{五门核心考试对比}\label{tab:12-exams}
\tablefont
\begin{tabular}{@{}L{2.0cm}L{3.0cm}C{1.5cm}C{1.6cm}C{1.5cm}L{2.6cm}@{}}
\toprule
\textbf{考试} & \textbf{定位} & \textbf{计分题} & \textbf{时长} & \textbf{通过线} & \textbf{前置/担保} \\
\midrule
SIE & 行业基础（人人可考） & 75 & 1h45m & 70\% & 无前置、无需担保 \\
Series 7 & 通用证券代表 & 125 & 3h45m & 72\% & 需 SIE + 雇主担保 \\
Series 57 & 证券交易员 & 50 & 1h45m & 70\% & 需 SIE + 雇主担保 \\
Series 65 & 投资顾问法 & 130 & 3h & 72\% & 无担保要求 \\
Series 24 & 通用证券负责人 & 150 & 3h45m & 70\% & 需 Series 7 类资格 + 担保 \\
\bottomrule
\end{tabular}
\end{table}

\begin{keybox}
这五门考试的分工可以这样理解：\textbf{SIE 考「行业常识」，Series 7 考「产品与客户」，
Series 57 考「交易台规则」，Series 65 考「受托责任与投资建议」，Series 24 考「监督与合规管理」}。
对量化从业者而言，Series 57 与 Series 24 的内容最接近日常工作场景（做市义务、订单报告、净资本）。
\end{keybox}

\section{SIE：证券行业基础考试}

\subsection{考试概述与知识领域}

\begin{table}[htbp]
\centering
\caption{SIE 四大知识领域及权重}\label{tab:12-sie}
\tablefont
\begin{tabularx}{\textwidth}{@{}C{1.2cm}L{4.4cm}Y@{}}
\toprule
\textbf{权重} & \textbf{领域} & \textbf{核心考点} \\
\midrule
16\% & 资本市场 & SEC/FINRA/MSRB/美联储职能、一级与二级市场、ATS 与暗池、货币政策与财政政策 \\
44\% & 产品理解 & 权益（普通股/优先股/ADR/认股权证）、债券（国债/市政债/公司债）、打包产品（基金/ETF/REIT/年金）、期权基础、风险类型 \\
31\% & 交易、账户与禁止行为 & 现金账户与保证金账户、订单类型、T+1 结算、KYC 与适合性、内幕交易（10b-5）、AML \\
9\% & 监管框架 & 1933 证券法、1934 证券交易法、1940 投资公司法、1940 投资顾问法 \\
\bottomrule
\end{tabularx}
\end{table}

\subsection{核心计算：债券收益率与税收等价}

SIE 与 Series 7 都会考债券的三种收益率。设年票息 $C$、面值 $F$、当前价格 $P$、剩余期数 $n$：

\begin{align}
  \text{当前收益率：}\quad & CY = \frac{C}{P} \label{eq:12-cy}\\
  \text{到期收益率：}\quad & P = \sum_{t=1}^{n}\frac{C}{(1+y)^{t}}+\frac{F}{(1+y)^{n}} \label{eq:12-ytm}\\
  \text{税收等价收益率：}\quad & TEY = \frac{r_{\text{muni}}}{1-t} \label{eq:12-tey}
\end{align}

其中 \cref{eq:12-tey} 是市政债与应税债券的比较工具：$t$ 为投资者的边际税率，
$r_{\text{muni}}$ 为免税市政债收益率。

\begin{example}[市政债与公司债的选择]\label{ex:12-tey}
某投资者边际税率 $t = 35\%$（联邦 + 州）。免税市政债收益率 $3.2\%$，
一只同等风险的公司债收益率 $4.6\%$。
市政债的税收等价收益率为 $TEY = 0.032/(1-0.35) = 4.92\%$，
\textbf{高于公司债的 $4.6\%$，因此市政债更优}。
换个角度看：公司债的税后收益为 $4.6\%\times(1-0.35) = 2.99\% < 3.2\%$，结论一致。
\textbf{当税率较低（例如 $t=12\%$）时，$TEY = 3.64\% < 4.6\%$，应税债券反而更优}——
这就是为什么市政债的主要持有者是高税率个人与保险公司。
\end{example}

\begin{figure}[htbp]
\centering
\begin{tikzpicture}
\begin{axis}[
  width=11.5cm, height=6.0cm, grid=major,
  xlabel={边际税率 $t$}, ylabel={税收等价收益率（\%）},
  domain=0:0.5, samples=100,
  legend pos=north west, legend style={font=\scriptsize},
  ymin=2.5, ymax=9, xtick={0,0.1,0.2,0.3,0.4,0.5}
]
\addplot[quantblue, very thick] {3.2/(1-x)};
\addlegendentry{市政债 3.2\%}
\addplot[quantgreen, very thick, dashed] {4.0/(1-x)};
\addlegendentry{市政债 4.0\%}
\addplot[quantred, thick, dotted] {4.6};
\addlegendentry{应税公司债 4.6\%（水平线）}
\end{axis}
\end{tikzpicture}
\caption{税收等价收益率随边际税率的变化：与应税债券水平线的交点即为「无差异税率」}\label{fig:12-tey}
\end{figure}

\subsection{核心计算：期权基础盈亏}

\begin{table}[htbp]
\centering
\caption{四个基本期权头寸的盈亏平衡与最大盈亏}\label{tab:12-options}
\tablefont
\begin{tabular}{@{}L{2.6cm}L{3.0cm}L{3.0cm}L{4.0cm}@{}}
\toprule
\textbf{头寸} & \textbf{盈亏平衡点} & \textbf{最大盈利} & \textbf{最大亏损} \\
\midrule
买入看涨 & $K+\text{权利金}$ & 无上限 & 权利金 \\
卖出看涨 & $K+\text{权利金}$ & 权利金 & 无上限 \\
买入看跌 & $K-\text{权利金}$ & $K-\text{权利金}$ & 权利金 \\
卖出看跌 & $K-\text{权利金}$ & 权利金 & $K-\text{权利金}$ \\
\midrule
备兑看涨 & $S_0-\text{权利金}$ & $(K-S_0)+\text{权利金}$ & $S_0-\text{权利金}$ \\
保护性看跌 & $S_0+\text{权利金}$ & 无上限 & $S_0-\text{权利金}$ \\
\bottomrule
\end{tabular}
\end{table}

\begin{exambox}[Series 7 高频考点]
期权在 Series 7 中约占 30--35 题，几乎全部围绕三件事：
(1) \textbf{画盈亏图（Profit/Loss Diagram）}；(2) \textbf{算盈亏平衡点}；(3) \textbf{算保证金要求}。
考试技巧：所有计算都从「最大风险 = 最大损失」出发，
因为最大损失决定了保证金要求，而保证金计算是失分最多的题型。
\end{exambox}

\begin{figure}[htbp]
\centering
\begin{tikzpicture}[scale=0.9, every node/.style={font=\small\sffamily}]
  \draw[->, thick, inkgray] (0,0) -- (8.6,0) node[right] {股价 $S_T$};
  \draw[->, thick, inkgray] (0,-1.5) -- (0,2.2) node[above] {损益};
  % 备兑看涨
  \draw[quantblue, very thick] (0,-1.0) -- (3.0,-1.0) -- (7.6,1.2);
  \draw[dashed, inkgray!55] (3.0,-1.5) -- (3.0,2.0);
  \draw[dashed, inkgray!55] (1.8,-1.5) -- (1.8,2.0);
  \node[below, inkgray, font=\scriptsize] at (1.8,-1.55) {盈亏平衡 $48$};
  \node[below, inkgray, font=\scriptsize] at (3.0,-1.55) {行权价 $55$};
  \node[quantblue, above, font=\scriptsize] at (6.0,1.05) {备兑看涨（上限）};
  \node[quantblue, below, font=\scriptsize] at (0.9,-0.95) {买入价 $50$};
  \draw[inkgray!40, dotted] (0.6,0) -- (8.2,0);
\end{tikzpicture}
\caption{备兑看涨的损益图（买入 50 元股票、卖出 55 元看涨收权利金 2 元）：最大盈利 7 元、盈亏平衡 48 元}\label{fig:12-covered}
\end{figure}

\section{Series 7：通用证券代表}

\subsection{考试概述与知识领域}

Series 7 是注册代表的旗舰考试，通过后可交易几乎所有证券产品。知识领域可归纳为四块：

\begin{enumerate}[leftmargin=1.8em]
  \item \textbf{为客户寻求业务}：客户画像、税务考量（短期与长期资本利得、成本基础计算 FIFO/LIFO/特定识别）、
        退休计划（传统/Roth IRA、401(k)、403(b)、SEP、SIMPLE）、教育储蓄（529、Coverdell ESA）；
  \item \textbf{评估客户财务画像与风险}：各类证券产品、市政债（GO 与 Revenue）、
        美国政府证券（T-Bills/Notes/Bonds、TIPS、STRIPS）、公司债、打包产品、期权策略、保证金账户；
  \item \textbf{开立、维护与关闭账户}：账户文档、确认与对账单、投诉处理与记录保存；
  \item \textbf{获取并验证客户买卖指示}：订单类型、执行与结算、交易报告。
\end{enumerate}

\subsection{核心计算：保证金账户与追缴}

美国证券交易的保证金规则由\textbf{Regulation T}（初始保证金，通常 50\%）与
\textbf{FINRA 4210}（维持保证金，多头通常 25\%）共同约束。

以多头为例，设买入价 $P_0$、股数 $N$、初始保证金比例 $m_0$、维持保证金比例 $m$：
借款额 $L = N P_0(1-m_0)$，账户权益为 $E = NP - L$。\textbf{追缴（margin call）发生在
权益恰好等于维持保证金要求时}：

\begin{equation}
  \frac{NP_{\text{call}}-L}{NP_{\text{call}}} = m
  \quad\Longrightarrow\quad
  P_{\text{call}} = \frac{L}{N(1-m)} = \frac{P_0(1-m_0)}{1-m}
  \label{eq:12-call}
\end{equation}

\begin{example}[多头保证金追缴价的计算]\label{ex:12-longcall}
以每股 60 元买入 100 股，共 6000 元，按 Reg T 借入 50\% 即 3000 元，维持保证金 25\%。
由 \cref{eq:12-call}：
\[
  P_{\text{call}} = \frac{60\times(1-0.5)}{1-0.25} = \frac{30}{0.75} = 40 \text{ 元}.
\]
验证：价格跌到 40 元时，市值 4000 元、权益 $4000-3000 = 1000$ 元，
恰好为市值的 25\%。\textbf{若价格继续下跌，投资者必须在规定时间内补足保证金，
否则经纪商将强制平仓}。注意初始保证金越低，追缴线越高（离现价越近），
这就是「高杠杆容易被震出」的数学原因。
\end{example}

\begin{figure}[htbp]
\centering
\begin{tikzpicture}
\begin{axis}[
  width=11.5cm, height=6.0cm, grid=major,
  xlabel={股价（元）}, ylabel={金额（元）},
  domain=20:80, samples=60,
  legend pos=north west, legend style={font=\scriptsize},
  ymin=-3000, ymax=6000
]
\addplot[quantblue, very thick] {100*x - 3000};
\addlegendentry{账户权益 $E = 100P - 3000$}
\addplot[quantred, very thick, dashed] {25*x};
\addlegendentry{维持保证金要求 $0.25\times100P$}
\addplot[inkgray, thick, dotted] {3000};
\addlegendentry{借款额（固定 3000 元）}
\addplot[only marks, mark=*, quantgold, mark size=2.4pt] coordinates {(40,1000)};
\node[quantgold, font=\scriptsize, anchor=south west] at (40,1000) {追缴点 $P=40$};
\end{axis}
\end{tikzpicture}
\caption{多头保证金账户：权益线与维持保证金要求的交点即追缴价}\label{fig:12-margin}
\end{figure}

\begin{notebox}[空头与期权的保证金差异]
空头的逻辑与多头相反：卖出股票的所得计入信用余额（credit balance），
维持保证金通常为 30\%，追缴价满足 $MV = \text{信用余额}/1.30$。
期权的保证金规则更复杂：\textbf{备兑看涨不需要额外保证金}（因为股票已作为担保），
而\textbf{裸卖看涨理论风险无限，其保证金要求通常为「权利金 + 标的市值的一定比例」}；
裸卖看跌的保证金则基于行权价与权利金。
考试中最容易被忽略的一点是：\textbf{价差策略（spread）的保证金只按两个行权价之差计算}，
而不是分别计算两条腿——这正是价差策略的资金效率来源。
\end{notebox}

\section{Series 57：证券交易员}

Series 57 面向股票、公司债与市政债交易员，是\textbf{与量化交易岗位最相关}的一门考试。

\begin{table}[htbp]
\centering
\caption{Series 57 的知识领域与交易台实务}\label{tab:12-57}
\tablefont
\begin{tabularx}{\textwidth}{@{}C{1.2cm}L{4.2cm}Y@{}}
\toprule
\textbf{权重} & \textbf{领域} & \textbf{核心考点} \\
\midrule
58\% & 交易活动与禁止行为 & 做市义务与报价规则、操纵行为（洗售 Wash Sales、Matched Orders、Painting the Tape、Spoofing/Layering）、MNPI 与内幕交易、最佳执行、禁止前端运行、净资本规则 15c3-1 \\
26\% & 订单处理与交易报告 & 订单审计追踪（OATS 迁移至 CAT）、大额交易报告、场外交易处理 \\
16\% & 一般证券知识与监管 & 监管机构分工、记录保存（17a-3/17a-4）、AML 与 OFAC 制裁合规 \\
\bottomrule
\end{tabularx}
\end{table}

\begin{exambox}[Spoofing 与 Layering：算法交易时代的红线]
\textbf{Spoofing} 指挂出并不打算成交的大额订单以影响市场对其他参与者判断，
在成交前撤单；\textbf{Layering} 则是在多个价位挂出多层虚假订单。
二者在 2010 年《多德--弗兰克法案》后被明确列为非法。
对量化团队的含义是：\textbf{任何「用订单流误导市场」的订单生成逻辑都不是策略优化，
而是违法行为}，必须在合规审查与代码层面同时设防（例如对「高频撤单率」设置硬性监控阈值）。
\end{exambox}

\section{Series 65：统一投资顾问法律考试}

Series 65 是成为注册投资顾问（RIA/IAR）的代表性资格，其核心是
\textbf{受托责任（Fiduciary Duty）}与\textbf{投资顾问法}。

\subsection{核心计算：绩效指标}

Series 65 会考察一组绩效与风险调整指标，务必区分分母：

\begin{align}
  \text{夏普比率：}\quad & SR = \frac{R_p-R_f}{\sigma_p} && \text{（分母为总风险）}\\
  \text{特雷诺比率：}\quad & TR = \frac{R_p-R_f}{\beta_p} && \text{（分母为系统性风险）}\\
  \text{詹森 Alpha：}\quad & \alpha_p = R_p-\left[R_f+\beta_p\left(R_m-R_f\right)\right] && \text{（CAPM 定价的偏离）}\\
  \text{信息比率：}\quad & IR = \frac{R_p-R_b}{\sigma\left(R_p-R_b\right)} && \text{（分母为跟踪误差）}
\end{align}

\begin{example}[同样 12\% 收益，哪个基金经理更好]\label{ex:12-perf}
无风险利率 3\%，市场收益 10\%。基金 A：$R_p = 12\%$、$\sigma_p = 18\%$、$\beta = 1.2$；
基金 B：$R_p = 12\%$、$\sigma_p = 12\%$、$\beta = 1.0$。则
\begin{align*}
  SR_A &= \frac{0.12-0.03}{0.18} = 0.50, &\quad
  TR_A &= \frac{0.09}{1.2} = 0.075, &\quad
  \alpha_A &= 0.006;\\
  SR_B &= \frac{0.09}{0.12} = 0.75, &\quad
  TR_B &= \frac{0.09}{1.0} = 0.09, &\quad
  \alpha_B &= 0.02.
\end{align*}
其中 $\alpha_A = 0.12-\left[0.03+1.2\times0.07\right] = 0.006$，
$\alpha_B = 0.12-\left[0.03+1.0\times0.07\right] = 0.02$。
\textbf{三个指标一致地给出 B 更优}：相同的绝对收益下，B 用更低的总风险与更低的系统性风险实现，
且超额收益中真正属于「选股能力」的部分（Alpha）更大。
考试要点是：\textbf{当组合充分分散时，特雷诺比率与夏普比率结论一致；
当组合未充分分散时，应当使用夏普比率，因为它惩罚了可分散风险}。
\end{example}

\subsection{税后与通胀调整收益}

\begin{align}
  \text{税后收益：}\quad & R_{\text{after-tax}} = R_{\text{pretax}}\left(1-t\right) \label{eq:12-aftertax}\\
  \text{实际（通胀调整）收益：}\quad & R_{\text{real}} = \frac{1+R_{\text{nominal}}}{1+i}-1 \approx R_{\text{nominal}}-i \label{eq:12-real}
\end{align}

\begin{exambox}[两个高频陷阱]
(1) \textbf{不能用「名义收益 − 通胀」当作精确的实际收益}，正确做法是
$\left(1+R\right)/\left(1+i\right)-1$（考试中的选项往往两者都给，需要用近似与精确区分）；
(2) \textbf{长期资本利得与短期资本利得的税率不同}，计算税后收益时必须先判断持有期，
且\textbf{市政债利息免征联邦所得税}——这也是 \cref{eq:12-tey} 存在的原因。
\end{exambox}

\section{Series 24：通用证券负责人}

Series 24 是管理级资格，考的是\textbf{监督与合规架构}而非产品细节。

\begin{table}[htbp]
\centering
\caption{Series 24 的四大领域}\label{tab:12-24}
\tablefont
\begin{tabularx}{\textwidth}{@{}C{1.2cm}L{4.4cm}Y@{}}
\toprule
\textbf{权重} & \textbf{领域} & \textbf{核心考点} \\
\midrule
24\% & 投资银行与研究监管 & 包销（Firm Commitment）与代销（Best Efforts）、注册流程与招股说明书、尽职调查、研究分析师隔离、安静期 \\
19\% & 交易与做市监管 & 做市商义务、Regulation SHO 卖空规则、大宗与异常交易报告 \\
29\% & 客户账户与销售实践 & 账户审批与监督、适合性与 Reg BI、AML 计划、隐私规则 Reg S-P、投诉处理、老年人财务剥削报告 \\
28\% & 运营与财务监管 & 净资本规则 15c3-1、客户保护规则 15c3-3、记录保存 17a-3/17a-4、FOCUS 报告、年度审计、业务连续性（BCP）与网络安全 \\
\bottomrule
\end{tabularx}
\end{table}

\begin{notebox}[净资本规则的两个直觉]
SEC Rule 15c3-1 要求在券商资产与负债之间维持一个最低的「净资本」缓冲，
其思想是：\textbf{保证券商在市场剧烈波动时仍有能力履行对客户的义务}。
Rule 15c3-3（客户保护规则）则要求把客户证券与自有证券隔离，
并通过「准备金账户（reserve bank account）」确保客户资产不被挪用。
对量化团队而言，这两条规则的实务含义是：\textbf{策略的资金占用不只包括保证金，
还包括券商的监管资本成本}——这也是 XVA 中 KVA（资本估值调整）存在的现实根源。
\end{notebox}

\section{考证路径与量化从业者的取舍}

\begin{table}[htbp]
\centering
\caption{典型考证路径对照}\label{tab:12-path}
\tablefont
\begin{tabular}{@{}L{3.0cm}L{3.6cm}L{4.0cm}@{}}
\toprule
\textbf{职业方向} & \textbf{必要考试} & \textbf{辅助/加分考试} \\
\midrule
注册代表 / 经纪人 & SIE + Series 7 & Series 63 或 66 \\
交易员（自营/做市） & SIE + Series 57 & Series 7 \\
投资顾问（RIA/IAR） & Series 65 & Series 7（可选） \\
合规 / 负责人 & Series 24 & Series 7、63/65 \\
量化研究员 / 分析师 & 通常无强制要求 & Series 7/57 有助于理解业务与合规边界 \\
\bottomrule
\end{tabular}
\end{table}

\begin{appbox}[量化从业者应该如何取舍]
对量化研究员而言，这些考试的价值不在于数学深度（考试的数学远简单于日常工作），而在于四点：
\begin{enumerate}[leftmargin=1.5em]
  \item \textbf{合规边界的可执行化}：知道哪些订单行为违法（Spoofing、前端运行、洗售），
        才能把合规规则写进代码与风控阈值；
  \item \textbf{业务语汇}：与交易台、销售、合规沟通时能使用共同术语，减少误解成本；
  \item \textbf{职业许可}：需要与客户互动或提供建议的岗位必须有相应资格；
  \item \textbf{产品广度}：考试覆盖了策略可能触及的几乎所有证券类型及其交易规则（结算周期、报告要求）。
\end{enumerate}
\end{appbox}

\begin{summarybox}
\begin{itemize}[leftmargin=1.4em]
  \item 五门考试的分工：SIE（行业常识）$\to$ Series 7（产品与客户）、57（交易台规则）、
        65（受托责任与投资建议）、24（监督与合规管理）。
  \item 债券的三种收益率：当前收益率 $C/P$、到期收益率（解 \cref{eq:12-ytm}）、
        税收等价收益率 $TEY = r_{\text{muni}}/(1-t)$。
  \item 多头保证金的追缴价 $P_{\text{call}} = P_0(1-m_0)/(1-m)$；
        杠杆越高，追缴线离现价越近。
  \item 期权的盈亏平衡与最大损失是保证金计算的基础；价差策略的保证金按行权价之差计算。
  \item 风险调整绩效指标的分母各不相同：夏普用 $\sigma_p$、特雷诺用 $\beta_p$、
        信息比率用跟踪误差；组合未充分分散时应看夏普比率。
  \item 税后收益 $R(1-t)$、实际收益 $(1+R)/(1+i)-1$；市政债利息免联邦税，这是 TEY 的用途。
  \item Series 24 的核心是净资本（15c3-1）与客户保护（15c3-3）——
        它们最终以资本成本的形式进入量化策略的成本核算。
\end{itemize}
\end{summarybox}

\begin{exercises}
\item 某投资者边际税率 24\%，免税市政债收益率 2.8\%。要使其与应税债券无差异，应税债券的收益率应为多少？
若市场上同等风险的应税债券报价 3.5\%，应选择哪一个？

\item 以每股 45 元买入 200 股，初始保证金 50\%，维持保证金 30\%。
求追缴价，并说明如果初始保证金降到 40\%，追缴价如何变化。

\item 某投资者买入 100 股（每股 50 元）并卖出 1 张 55 元看涨期权（权利金 2 元）。
请写出最大盈利、最大亏损与盈亏平衡点，并说明该策略在什么市场情形下最优。

\item 基金 C 的收益 15\%、标准差 20\%、$\beta = 1.3$；无风险利率 3\%，市场收益 10\%。
请计算夏普比率、特雷诺比率与詹森 Alpha，并判断其超额收益主要来自选股还是承担系统性风险。

\item 某投资者名义收益 8\%，通胀 3\%，适用税率 20\%，且收益全部为短期资本利得。
请计算其税后实际收益（分别用近似公式与精确公式）。

\item 解释为什么「裸卖看涨期权」的保证金要求高于「备兑看涨」，
并说明在 Series 24 的净资本规则视角下，这两种头寸对券商资本占用的差异。
\end{exercises}

\begin{furtherreading}
\refitem FINRA, \emph{Content Outlines for the SIE, Series 7, Series 57, Series 65 and Series 24 Examinations}, Financial Industry Regulatory Authority, 2024.（唯一的权威大纲，报考前必读）
\refitem U.S. Securities and Exchange Commission, \emph{Regulation Best Interest: The Broker-Dealer Standard of Conduct}, Release No. 34-86031, 2019.（适合性与最佳利益规则的法规原文）
\refitem Basel Committee on Banking Supervision, \emph{Basel III: Finalising Post-Crisis Reforms}, Bank for International Settlements, 2017.（与净资本、杠杆率要求对照阅读，理解资本约束的全貌）
\refitem Hull, J. C., \emph{Options, Futures, and Other Derivatives}, 10th ed., Pearson, 2018.（第 3 章关于保证金机制的讲解比考试教材更清晰）
\refitem Hasbrouck, J., \emph{Empirical Market Microstructure}, Oxford University Press, 2007.（理解做市义务、订单类型与交易报告背后的市场结构）
\end{furtherreading}
